\usepackage[a4paper, left=3.5cm,right=3.5cm,bottom=3.5cm]{geometry}
\usepackage[final]{graphicx}
\graphicspath{{../images/}}
\usepackage{mathrsfs}
\usepackage{newtxtext}
\usepackage{float}
\usepackage{comment}
\usepackage{tikz}
\usetikzlibrary{calc}
\usepackage{eso-pic}
\usepackage[dutch]{babel}
\usepackage{fancyhdr}
\usepackage{multirow}
\usepackage[table]{xcolor}
\usepackage[hidelinks]{hyperref}
\usepackage[toc, acronym, nomain]{glossaries-extra}
\usepackage{attachfile}
\usepackage{array}
\usepackage{float}

% Mooi hoekje links bovenin net als in de Voortman template
\newcounter{pagecount}
\AtBeginShipout{%
	\stepcounter{pagecount} % Tel de pagina's
	\ifnum\value{pagecount}>1 % Pas toe vanaf pagina 2
	\AtBeginShipoutUpperLeft{%
		\begin{tikzpicture}[remember picture, overlay]
			\fill[customred] (current page.north west) ++(3.0cm,0) -- (current page.north west) -- ++(0,-6.5cm) -- cycle;
		\end{tikzpicture}
	}
	\fi
}

\makeglossaries

\pagestyle{fancy}

\fancyhf{}

\definecolor{customred}{HTML}{D78C91}

\newenvironment{myfont}{\fontfamily{phv}\selectfont}{\par}

\fancyhead[L]{\textbf{Afstudeerstage}}
\fancyhead[C]{}
\fancyhead[R]{\today}

% Voettekst-instellingen
\fancyfoot[L]{Auteur: Florent Kegler}
\fancyfoot[C]{Pagina \thepage}
\fancyfoot[R]{\includegraphics[width=4cm]{VoortmanLogo}}

% maak de tabel ietsje groter want de standaard is soms beetje smal
\setlength{\arrayrulewidth}{0.5mm}
\setlength{\tabcolsep}{18pt}
\renewcommand{\arraystretch}{1.5}